% Why tagging lets OCaml move a live block and rewrite pointers,
% while C cannot. Before / after. Used in M10-L02.
\documentclass[tikz,border=4pt]{standalone}
\usepackage{tikz}
\input{_m10-styles.texinput}
\begin{document}
\begin{tikzpicture}[m10, x=1mm, y=-1mm, font=\large]
  % ---- BEFORE ----
  \node[note, font=\Large\bfseries, anchor=east] at (13,2) {before:};
  \node[stackcell, minimum width=18mm, minimum height=12mm, font=\large] (r1) at (24,2) {root};
  \node[freedcell, hatched, minimum width=16mm, minimum height=12mm] (g1) at (48,2) {};
  \node[heapcell, minimum width=18mm, minimum height=12mm, font=\large] (live1) at (70,2) {live};
  \draw[flow, line width=1pt] (r1.east) -- (live1.west);
  \node[badnote, anchor=west, font=\large] at (82,2) {garbage gap};

  % ---- AFTER ----
  \node[note, font=\Large\bfseries, anchor=east] at (13,26) {after:};
  \node[stackcell, minimum width=18mm, minimum height=12mm, font=\large] (r2) at (24,26) {root};
  \node[heapcell, minimum width=18mm, minimum height=12mm, font=\large] (live2) at (48,26) {live};
  \draw[goodflow, line width=1pt] (r2.east) -- (live2.west);
  \node[goodnote, anchor=west, font=\large] at (60,26) {moved, pointer rewritten};

  % the rule that makes it possible
  \node[safebox, minimum width=120mm, align=left, anchor=north west, font=\large]
    (rule) at (0,44)
    {\textbf{OCaml:} every word is tagged, so the GC knows which are
       pointers and can rewrite them. Moving is safe.};
  \node[warnbox, minimum width=120mm, align=left, anchor=north west, font=\large]
    (crule) at (0,62)
    {\textbf{C:} a word may be an \texttt{int} or a \texttt{void *};
       the runtime cannot tell, so it must never move a block.};
\end{tikzpicture}
\end{document}
